\par \Th Степень минимального многочлена числа $\xi_n \in \Comp$ над $\Q$ равна $\varphi(n)$.

\par \Proof Для произвольного $n \in \N$ положим 

$$\psi_n(x):=\prod_{\substack{1\leq d \leq n, \\ (d,n)=1}} (x-\xi_n^d)$$

\par Докажем, что это минимальный многочлен для $\xi_n$ над $\Q$. Проверим сначала, что $\psi_n \in \Z[x]$.
Действительно, поскольку выполнено равенство $\prod_{k|n} \psi_k(x)=x^n-1$ (???), то целочисленность всех коэффициентов многочлена $\psi_n$ легко проверяется по индукции.

\par Пусть теперь $m \in \Q[x]$ — минимальный многочлен для $\xi_n$ (возможно, отличный от $\psi_n$), $\xi$ — его произвольный корень. Покажем, что $m$ должен также иметь корни $\xi^p$ для всех
простых чисел $p \nmid n$. В силу минимальности, $m | x^n-1$, то есть $x^n-1=m(x)g(x)$ для некоторого $g \in \Q[x]$, причем коэффициенты многочленов $m,g$ целочисленны, поскольку многочлен $x^n-1$ примитивен и имеет старший коэффициент, равный 1. 

\par Предположим, что $m(\xi_p)\neq 0$, тогда $g(\xi_p)=0$, поэтому $m | g(x^p)$, то есть выполнено равенство $g(x^p)=m(x)h(x)$ для некоторого $h \in \Q[x]$, и, аналогично, коэффициенты многочлена $h$ целочисленны. В поле $\Z_p$ равенство принимает вид $\overline{g}(x^p)=\overline{g}(x)^p=\overline{m}(x)\overline{h}(x)$. Значит, в алгебраическом замыкании поля $\Z_p$ многочлены $\overline{g}$ и $\overline{m}$ имеют общий корень. Но многочлен $x^n-1$ не имеет кратных корней, поскольку его формальная производная равна $nx^{n-1}$ и не имеет корней, отличных от нуля. Значит, $m(\xi^p)=0$. Но тогда и для всех чисел $d \in \N$ таких, что $(d,n)=1$, выполнено $m(\xi_n^d)=0$. Значит, минимальный многочлен $m$ совпадает с $\psi_n$ и имеет степень, равную $\varphi(n)$. \EndProof

\par \Cor Пусть $n \in \N$. Тогда правильный $n$-угольник можно построить циркулем и линейкой $\Leftrightarrow \varphi(n)=2^k, k\in \N$
\par \Proof Очевидное следствие из этой теоремы и теоремы из билета 14 на отл \EndProof